%----------------------------------------------------------------------------------------
% Preambulo y Configuración
%----------------------------------------------------------------------------------------

\documentclass[
    11pt,
    spanish,
    singlespacing,
    parskip,
    headsepline,
    bookmarks=true,
    unicode=true,
    pdftoolbar=true,
    pdfmenubar=true,
    pdffitwindow=false,
    colorlinks=true,
    linkcolor=blue,
    citecolor=blue,
    urlcolor=blue
]{MastersDoctoralThesis}

\usepackage[utf8]{inputenc} % Codificación de entrada UTF-8
\usepackage[T1]{fontenc}    % Codificación de salida para caracteres especiales
\usepackage{graphicx}       % Manejo de gráficos
\usepackage{eso-pic}        % Permite agregar fondos
\usepackage{hyperref}       % Manejo de hipervínculos y marcadores

% Redefinición de caracteres problemáticos en marcadores
\hypersetup{
    pdftitle={Sistema de telemetría y detección de eventos para motocicletas},
    pdfauthor={Tomás Enrique Botalla},
    pdfkeywords={Sistemas embebidos, telemetría, motocicletas, detección de eventos},
    pdfstartview={FitH},
    unicode=true,
    colorlinks=true,
    linkcolor=blue,
    citecolor=blue,
    urlcolor=blue
}

\pdfstringdefDisableCommands{%
  \def\texttt#1{#1}%
  \def\textbf#1{#1}%
  \def\textit#1{#1}%
  \def\"{\"}%
  \def\~{~}%
  \def\'{'}%
  \def\^{}%
  \def\textunderscore{\_} % Manejo del subrayado en marcadores
}


% Definir comandos requeridos por la clase
\newcommand{\degreename}{Carrera de Especialización en Sistemas Embebidos}
\newcommand{\univname}{Universidad de Buenos Aires}
\newcommand{\keywordnames}{Palabras clave:}
%----------------------------------------------------------------------------------------
% Documento Principal
%----------------------------------------------------------------------------------------

\begin{document}

% Configuración de la portada
\posgrado{Carrera de Especialización en Sistemas Embebidos}
\keywords{Sistemas embebidos, telemetría, motocicletas, detección de eventos}

% Incluir la portada desde un archivo separado
%----------------------------------------------------------------------------------------
% PORTADA
%----------------------------------------------------------------------------------------
\begin{titlepage}
    % Fondo completo con el PDF que incluye la barra y el logo
    \AddToShipoutPictureBG*{\includegraphics[width=\paperwidth, height=\paperheight]{Figures/fondo.pdf}}

    % Contenido principal
    \begin{flushright}
        \setlength{\rightskip}{-2cm} % Ajusta la sangría derecha
        \vspace*{7.5cm} % Ajustar según la posición vertical deseada

        % Título
        {\fontfamily{phv}\bfseries\fontsize{33pt}{40pt}\selectfont
        Sistema de telemetría y detección de eventos para motocicletas} \\[1.5cm]

        % Autor
        {\fontfamily{phv}\fontsize{20pt}{25pt}\selectfont
        Ing. Tomás Enrique Botalla} \\[1cm]

        % Carrera o Maestría (comentar o descomentar la línea correspondiente)
        {\fontfamily{phv}\fontsize{15pt}{20pt}\selectfont
        \textbf{Carrera de Especialización en Sistemas Embebidos}
        % \textbf{Carrera de Especialización en Internet de las Cosas} \\
        % \textbf{Carrera de Especialización en Inteligencia Artificial} \\
        % \textbf{Maestría en Sistemas Embebidos} \\
        % \textbf{Maestría en Internet de las Cosas} \\
        % \textbf{Maestría en Inteligencia Artificial Embebida} \\
        % \textbf{Maestría en Computación de Borde} \\
        % \textbf{Maestría en Inteligencia Artificial} \\
        } \\[2cm]

        % Director
        {\fontfamily{phv}\fontsize{11pt}{15pt}\selectfont
        \textbf{Director:} Esp. Ing. Germán Eloy Cardozo (UNLaM)} \\[1cm]

        % Jurados
        {\fontfamily{phv}\fontsize{11pt}{15pt}\selectfont
        \textbf{Jurados:}} \\[0.5cm]
        {\fontfamily{phv}\fontsize{11pt}{15pt}\selectfont
        Por definir} \\ 
        {\fontfamily{phv}\fontsize{11pt}{15pt}\selectfont
        Por definir} \\ 
        {\fontfamily{phv}\fontsize{11pt}{15pt}\selectfont
        Por definir} \\[2cm]

        % Fecha y lugar
        {\fontfamily{phv}\itshape\fontsize{10pt}{12pt}\selectfont
        Ciudad Autónoma de Buenos Aires, Noviembre de 2026}
    \end{flushright}
\end{titlepage}


% Configuración del contenido preliminar
\frontmatter % Usar numeración romana para las páginas preliminares
\pagestyle{plain} % Estilo de encabezado simple

%----------------------------------------------------------------------------------------
% Resumen
%----------------------------------------------------------------------------------------

\begin{abstract}
\addchaptertocentry{\abstractname}
En la presente memoria se describe el diseño e implementación de un sistema embebido de telemetría y detección de eventos para motocicletas. El trabajo integra adquisición de variables cinemáticas y geográficas, almacenamiento en tarjeta de memoria removible y comunicación con una aplicación móvil en el teléfono.

Para su implementación fueron clave conocimientos relacionados con programación de microcontroladores, protocolos de comunicación, pruebas de software, sistemas operativos de tiempo real y gestión de proyectos.
\end{abstract}

%----------------------------------------------------------------------------------------
% Agradecimientos
%----------------------------------------------------------------------------------------

% \begin{acknowledgements}
% \vspace{1.5cm}
% Esta sección es para agradecimientos personales y es totalmente \textbf{OPCIONAL}.
% \end{acknowledgements}

%----------------------------------------------------------------------------------------
% Índice
%----------------------------------------------------------------------------------------

\tableofcontents
\listoffigures
\listoftables

%----------------------------------------------------------------------------------------
% Dedicatoria
%----------------------------------------------------------------------------------------

% \dedicatory{\textbf{Dedicado a... [OPCIONAL]}}

%----------------------------------------------------------------------------------------
% Capítulos
%----------------------------------------------------------------------------------------

\mainmatter % Iniciar numeración numérica para el contenido principal
\pagestyle{thesis} % Estilo de encabezado de tesis

% Incluir capítulos desde archivos separados
% Chapter 1

\chapter{Introducción general}

\label{Chapter1}
\label{IntroGeneral}

En este capítulo se introduce el dominio de la telemetría embebida aplicada a motocicletas. Se presentan el contexto, la motivación, el estado del arte y los objetivos del trabajo.

%----------------------------------------------------------------------------------------

% Define some commands to keep the formatting separated from the content
\newcommand{\keyword}[1]{\textbf{#1}}
\newcommand{\tabhead}[1]{\textbf{#1}}
\newcommand{\code}[1]{\texttt{#1}}
\newcommand{\file}[1]{\texttt{\bfseries#1}}
\newcommand{\option}[1]{\texttt{\itshape#1}}
\newcommand{\grados}{$^{\circ}$}

%----------------------------------------------------------------------------------------
\section{Introducción general al tema}
\label{sec:intro-general-tema}

Existe una amplia variedad de sistemas embebidos en la actualidad. Se encuentran en lo cotidiano en electrodomésticos, equipos médicos, instrumentación industrial, sistemas de transporte, entre muchos otros. En todos los casos, un firmware dedicado ejecuta la lógica que define la funcionalidad del producto y condiciona su confiabilidad en el tiempo.

Un sistema embebido es un sistema computacional orientado a una función acotada, integrado en un producto mayor y sujeto a límites de tamaño, consumo, costo y tiempo de respuesta. A diferencia de un equipo de uso general, el hardware y el software se conciben de forma conjunta para adquirir señales, controlar un proceso o registrar información en tiempo acotado.

En el ámbito vehicular, los sistemas embebidos abordan distintos dominios de aplicación. Uno de ellos es la telemetría vehicular, que concentra la adquisición, sincronización, almacenamiento y transmisión de magnitudes cinemáticas, y físicas. En competición y en flotas profesionales, esas soluciones están ampliamente extendidas. En el uso particular del conductor, la demanda se orienta hacia equipos compactos, confiables y de costo moderado.

%----------------------------------------------------------------------------------------
\section{Contexto y motivación}
\label{sec:contexto-motivacion}

La conducción en motocicleta impone condiciones dinámicas exigentes. El usuario enfrenta cambios rápidos de velocidad, inclinación en curvas y mayor exposición a maniobras bruscas o a incidentes. Un dispositivo embebido en el vehículo puede actuar como instrumento de medición y como respaldo de datos, en la línea del concepto de ``caja negra'' usado en otros medios de transporte.

El proceso que se busca instrumentar es ese entorno de conducción. Se requiere registrar variables cinemáticas y geográficas, identificar maniobras bruscas o situaciones de riesgo y conservar un historial utilizable después del recorrido o ante un incidente. Muchos conductores recurren a aplicaciones en el teléfono móvil. Esas aplicaciones no están integradas al vehículo, dependen de sensores del teléfono y no garantizan registro autónomo en el vehículo. Los dispositivos profesionales ofrecen alta confiabilidad, pero exigen instalación, configuración y mantenimiento más complejos, además de costos elevados y ecosistemas cerrados poco accesibles para un uso particular.

Las restricciones de tamaño, consumo y costo se combinan con condiciones físicas exigentes en el vehículo, como vibración, variación térmica y espacio limitado para el montaje. Esas limitaciones obligan a una arquitectura optimizada. El microcontrolador debe coordinar periféricos en paralelo, sostener tasas de muestreo adecuadas y ejecutar lógica de análisis en tiempo acotado sin depender de servicios remotos permanentes. La plataforma debe reunir adquisición multivariable, persistencia local y enlace de corto alcance hacia una aplicación móvil en el teléfono.

La motivación del trabajo es cubrir esa brecha con un sistema embebido dedicado. Se prioriza la medición orientada al vehículo, el almacenamiento en tarjeta de memoria removible y el enlace inalámbrico sin conectividad celular en el dispositivo. El prototipo integra un microcontrolador \code{STM32} (STMicroelectronics) sobre la placa de desarrollo \code{NUCLEO-F446RE}, sensado inercial, receptor de posicionamiento global, registro en tarjeta microSD y comunicación por Bluetooth de baja energía (BLE) con la aplicación móvil.

%----------------------------------------------------------------------------------------
\section{Estado del arte}
\label{sec:estado-arte}

Durante el relevamiento sobre telemetría embebida en motocicletas se constató que el mercado profesional está dominado por plataformas cerradas, probadas en competición y en desarrollo de producto \citep{garcia2026_moto_telemetry}. Fabricantes como AiM Technologies, MoTeC y 2D Datarecording ofrecen dispositivos con unidad de medición inercial integrada, posicionamiento global, interfaces de bus de área de controlador (CAN) y software propietario de análisis \citep{aim_evo5,aim_mxl2,motec_c187,data2d_moto}. Esas soluciones alcanzan alta confiabilidad, pero suelen implicar costos elevados, formatos de datos poco abiertos y poca extensibilidad para algoritmos propios de análisis de conducción.

En la tabla~\ref{tab:comercial-telemetria} se resume un conjunto de las alternativas comerciales de nivel profesional relevadas.

\begin{table}[htbp]
	\centering
	\caption{Alternativas comerciales representativas de telemetría para vehículos y motocicletas.}
	\label{tab:comercial-telemetria}
	\begin{tabular}{|p{2.2cm}|p{2.4cm}|p{4.2cm}|p{4.5cm}|}
		\hline
		\tabhead{Producto} & \tabhead{Fabricante} & \tabhead{Capacidades principales} & \tabhead{Limitaciones} \\
		\hline
		EVO5 & AiM Technologies \citep{aim_evo5} & Canales configurables, posicionamiento global, bus CAN, software Race Studio & Arquitectura cerrada. Costo elevado en configuraciones completas \\
		\hline
		MXL2 & AiM Technologies \citep{aim_mxl2} & Acelerómetro y giroscopio internos, posicionamiento global, entradas analógicas, memoria interna & Mismo ecosistema propietario \\
		\hline
		C187 & MoTeC \citep{motec_c187} & Display a color, registro de datos, posicionamiento global, software MoTeC i2 & Enfoque competición. Hardware y licencias de alto costo \\
		\hline
		Sticklogger V2/V3/V4 & 2D Datarecording \citep{data2d_moto} & Registro en memoria USB 3.0, señales CAN, analógicas y digitales, configuración vía USB & Plataforma cerrada. Enfoque competición y entornos exigentes \\
		\hline
	\end{tabular}
\end{table}

En el ámbito académico, las publicaciones abordan registradores de eventos, reconstrucción de trayectorias, detección de colisiones y plataformas embebidas de bajo costo. Gelmini y col. propusieron un algoritmo de detección y clasificación de colisiones para vehículos de dos ruedas a partir de mediciones inerciales y de posicionamiento global \citep{gelmini2021_crash}. Smaiah y col. reconstruyeron trayectorias de motocicleta con un sistema multisenorial de bajo costo \citep{smaiah2018_trajectory}.

García Cuadra y col. publicaron una plataforma embebida abierta sobre microcontrolador STM32H7, con posicionamiento global de alta precisión, unidad de medición inercial, bus CAN con extensión CAN-FD y enlace celular \citep{garcia2026_moto_telemetry}. Ese aporte evidencia la viabilidad de arquitecturas abiertas frente a registradores cerrados de competición.

%----------------------------------------------------------------------------------------
\section{Objetivos del trabajo}
\label{sec:objetivos}

\subsection{Objetivo general}
\label{sec:objetivo-general}

El objetivo general del trabajo fue diseñar e implementar un sistema embebido de telemetría y detección de eventos para motocicletas, capaz de adquirir variables cinemáticas y geográficas, registrar datos de forma local y comunicarse con una aplicación móvil en el teléfono.

\subsection{Objetivos específicos}
\label{sec:objetivos-especificos}

Los objetivos específicos fueron los siguientes:

\begin{enumerate}
	\item Seleccionar e integrar una plataforma de hardware basada en la familia de microcontroladores \code{STM32} y los módulos de sensado necesarios para posicionamiento global y medición inercial.
	\item Desarrollar el firmware que permita adquirir, sincronizar y registrar muestras de sensores con tasas adecuadas para el análisis de la conducción.
	\item Implementar almacenamiento persistente en una tarjeta de memoria removible, con formatos de registro que permitan un análisis posterior.
	\item Establecer comunicación BLE con la aplicación móvil para la transferencia de telemetría.
	\item Incorporar lógica para identificar maniobras de conducción relevantes a partir de las señales adquiridas, como base para puntuación y alertas.
	\item Verificar el funcionamiento del sistema mediante ensayos sobre el hardware, el firmware y la integración con la aplicación móvil.
\end{enumerate}

\subsection{Alcance}
\label{sec:alcance}

El alcance comprendió el diseño e implementación de un prototipo embarcado sobre placa \code{NUCLEO-F446RE}, el firmware asociado y la aplicación móvil para el envío de telemetría por BLE. Se priorizó la adquisición multivariable, la persistencia local y la interfaz con el usuario móvil.

En el firmware se incluyó el uso de un sistema operativo de tiempo real para organizar tareas concurrentes: muestreo de sensores, registro en tarjeta de memoria removible y enlace BLE. Con ello se buscó separar las rutinas periódicas de las operaciones de entrada-salida de mayor latencia y mantener tiempos de respuesta acotados en el dispositivo.

\begin{figure}[htbp]
	\centering
	\includegraphics[width=0.9\textwidth]{"Figures/Diagrama bloques.png"}
	\caption{Diagrama de bloques del dispositivo.}
	\label{fig:alcance-diagrama-bloques}
\end{figure}

Quedaron fuera del alcance cualquier tipo de certificación u homologación del dispositivo, la conexión con la electrónica de la motocicleta mediante bus CAN u otros enlaces de diagnóstico y el uso de servicios en la nube para almacenamiento remoto. Tampoco se buscó igualar el desempeño ni el costo de dispositivos comerciales de competición. El foco fue proveer una solución extensible y de costo moderado para conductores particulares.

\chapter{Introducción específica} % Main chapter title

\label{Chapter2}

%----------------------------------------------------------------------------------------
%	SECTION 1
%----------------------------------------------------------------------------------------
Todos los capítulos deben comenzar con un breve párrafo introductorio que indique cuál es el contenido que se encontrará al leerlo.  La redacción sobre el contenido de la memoria debe hacerse en presente y todo lo referido al proyecto en pasado, siempre de modo impersonal.

\section{Estilo y convenciones}
\label{sec:ejemplo}

\subsection{Uso de mayúscula inicial para los título de secciones}

Si en el texto se hace alusión a diferentes partes del trabajo referirse a ellas como capítulo, sección o subsección según corresponda. Por ejemplo: ``En el capítulo \ref{Chapter1} se explica tal cosa'', o ``En la sección \ref{sec:ejemplo} se presenta lo que sea'', o ``En la subsección \ref{subsec:ejemplo} se discute otra cosa''.

Cuando se quiere poner una lista tabulada, se hace así:

\begin{itemize}
	\item Este es el primer elemento de la lista.
	\item Este es el segundo elemento de la lista.
\end{itemize}

Notar el uso de las mayúsculas y el punto al final de cada elemento.

Si se desea poner una lista numerada el formato es este:

\begin{enumerate}
	\item Este es el primer elemento de la lista.
	\item Este es el segundo elemento de la lista.
\end{enumerate}

Notar el uso de las mayúsculas y el punto al final de cada elemento.

\subsection{Este es el título de una subsección}
\label{subsec:ejemplo}

Se recomienda no utilizar \textbf{texto en negritas} en ningún párrafo, ni tampoco texto \underline{subrayado}. En cambio sí se debe utilizar \textit{texto en itálicas} para palabras en un idioma extranjero, al menos la primera vez que aparecen en el texto. En el caso de palabras que estamos inventando se deben utilizar ``comillas'', así como también para citas textuales. Por ejemplo, un \textit{digital filter} es una especie de ``selector'' que permite separar ciertos componentes armónicos en particular.

La escritura debe ser impersonal. Por ejemplo, no utilizar ``el diseño del firmware lo hice de acuerdo con tal principio'', sino ``el firmware fue diseñado utilizando tal principio''. 

El trabajo es algo que al momento de escribir la memoria se supone que ya está concluido, entonces todo lo que se refiera a hacer el trabajo se narra en tiempo pasado, porque es algo que ya ocurrió. Por ejemplo, "se diseñó el firmware empleando la técnica de test driven development".

En cambio, la memoria es algo que está vivo cada vez que el lector la lee. Por eso transcurre siempre en tiempo presente, como por ejemplo:

``En el presente capítulo se da una visión global sobre las distintas pruebas realizadas y los resultados obtenidos. Se explica el modo en que fueron llevados a cabo los test unitarios y las pruebas del sistema''.

Se recomienda no utilizar una sección de glosario sino colocar la descripción de las abreviaturas como parte del mismo cuerpo del texto. Por ejemplo, RTOS (\textit{Real Time Operating System}, Sistema Operativo de Tiempo Real) o en caso de considerarlo apropiado mediante notas a pie de página.

Si se desea indicar alguna página web utilizar el siguiente formato de referencias bibliográficas, dónde las referencias se detallan en la sección de bibliografía de la memoria, utilizado el formato establecido por IEEE.% \citep{IEEE:citation}. Por ejemplo, ``el presente trabajo se basa en la plataforma EDU-CIAA-NXP % \citep{CIAA}, la cual...''.

\subsection{Figuras} 

Al insertar figuras en la memoria se deben considerar determinadas pautas. Para empezar, usar siempre tipografía claramente legible. Luego, tener claro que \textbf{es incorrecto} escribir por ejemplo esto: ``El diseño elegido es un cuadrado, como se ve en la siguiente figura:''

\begin{figure}[h]
\centering
\includegraphics[scale=.45]{./Figures/cuadradoAzul.png}
\end{figure}

La forma correcta de utilizar una figura es con referencias cruzadas, por ejemplo: ``Se eligió utilizar un cuadrado azul para el logo, como puede observarse en la figura \ref{fig:cuadradoAzul}''.

\begin{figure}[ht]
	\centering
	\includegraphics[scale=.45]{./Figures/cuadradoAzul.png}
	\caption{Ilustración del cuadrado azul que se eligió para el diseño del logo.}
	\label{fig:cuadradoAzul}
\end{figure}

El texto de las figuras debe estar siempre en español, excepto que se decida reproducir una figura original tomada de alguna referencia. En ese caso la referencia de la cual se tomó la figura debe ser indicada en el epígrafe de la figura e incluida como una nota al pie, como se ilustra en la figura \ref{fig:palabraIngles}.

\begin{figure}[htpb]
	\centering
	\includegraphics[scale=.3]{./Figures/word.jpeg}
	\caption{Imagen tomada de la página oficial del procesador\protect\footnotemark.}
	\label{fig:palabraIngles}
\end{figure}

\footnotetext{Imagen tomada de \url{https://goo.gl/images/i7C70w}}

La figura y el epígrafe deben conformar una unidad cuyo significado principal pueda ser comprendido por el lector sin necesidad de leer el cuerpo central de la memoria. Para eso es necesario que el epígrafe sea todo lo detallado que corresponda y si en la figura se utilizan abreviaturas entonces aclarar su significado en el epígrafe o en la misma figura.



\begin{figure}[ht]
	\centering
	\includegraphics[scale=.37]{./Figures/questionMark.png}
	\caption{¿Por qué de pronto aparece esta figura?}
	\label{fig:questionMark}
\end{figure}

Nunca colocar una figura en el documento antes de hacer la primera referencia a ella, como se ilustra con la figura \ref{fig:questionMark}, porque sino el lector no comprenderá por qué de pronto aparece la figura en el documento, lo que distraerá su atención.

Otra posibilidad es utilizar el entorno \textit{subfigure} para incluir más de una figura, como se puede ver en la figura \ref{fig:three graphs}. Notar que se pueden referenciar también las figuras internas individualmente de esta manera: \ref{fig:1de3}, \ref{fig:2de3} y \ref{fig:3de3}.
 
\begin{figure}[!htpb]
     \centering
     \begin{subfigure}[b]{0.3\textwidth}
         \centering
         \includegraphics[width=.65\textwidth]{./Figures/questionMark}
         \caption{Un caption.}
         \label{fig:1de3}
     \end{subfigure}
     \hfill
     \begin{subfigure}[b]{0.3\textwidth}
         \centering
         \includegraphics[width=.65\textwidth]{./Figures/questionMark}
         \caption{Otro.}
         \label{fig:2de3}
     \end{subfigure}
     \hfill
     \begin{subfigure}[b]{0.3\textwidth}
         \centering
         \includegraphics[width=.65\textwidth]{./Figures/questionMark}
         \caption{Y otro más.}
         \label{fig:3de3}
     \end{subfigure}
        \caption{Tres gráficos simples.}
        \label{fig:three graphs}
\end{figure}

El código para generar las imágenes se encuentra disponible para su reutilización en el archivo \file{Chapter2.tex}.

\subsection{Tablas}

Para las tablas utilizar el mismo formato que para las figuras, sólo que el epígrafe se debe colocar arriba de la tabla, como se ilustra en la tabla \ref{tab:peces}. Observar que sólo algunas filas van con líneas visibles y notar el uso de las negritas para los encabezados.  La referencia se logra utilizando el comando \verb|\ref{<label>}| donde label debe estar definida dentro del entorno de la tabla.

\begin{verbatim}
\begin{table}[h]
	\centering
	\caption[caption corto]{caption largo más descriptivo}
	\begin{tabular}{l c c}    
		\toprule
		\textbf{Especie}     & \textbf{Tamaño} & \textbf{Valor}\\
		\midrule
		Amphiprion Ocellaris & 10 cm           & \$ 6.000 \\		
		Hepatus Blue Tang    & 15 cm           & \$ 7.000 \\
		Zebrasoma Xanthurus  & 12 cm           & \$ 6.800 \\
		\bottomrule
		\hline
	\end{tabular}
	\label{tab:peces}
\end{table}
\end{verbatim}


\begin{table}[h]
	\centering
	\caption[caption corto]{caption largo más descriptivo.}
	\begin{tabular}{l c c}    
		\toprule
		\textbf{Especie} 	 & \textbf{Tamaño} 		& \textbf{Valor}  \\
		\midrule
		Amphiprion Ocellaris & 10 cm 				& \$ 6.000 \\		
		Hepatus Blue Tang	 & 15 cm				& \$ 7.000 \\
		Zebrasoma Xanthurus	 & 12 cm				& \$ 6.800 \\
		\bottomrule
		\hline
	\end{tabular}
	\label{tab:peces}
\end{table}

En cada capítulo se debe reiniciar el número de conteo de las figuras y las tablas, por ejemplo, figura 2.1 o tabla 2.1, pero no se debe reiniciar el conteo en cada sección. Por suerte la plantilla se encarga de esto por nosotros.

\subsection{Ecuaciones}
\label{sec:Ecuaciones}

Al insertar ecuaciones en la memoria dentro de un entorno \textit{equation}, éstas se numeran en forma automática  y se pueden referir al igual que como se hace con las figuras y tablas, por ejemplo ver la ecuación \ref{eq:metric}.

\begin{equation}
	\label{eq:metric}
	ds^2 = c^2 dt^2 \left( \frac{d\sigma^2}{1-k\sigma^2} + \sigma^2\left[ d\theta^2 + \sin^2\theta d\phi^2 \right] \right)
\end{equation}
                                                        
Para generar la ecuación \ref{eq:metric} se utilizó el siguiente código:

\begin{verbatim}
\begin{equation}
	\label{eq:metric}
	ds^2 = c^2 dt^2 \left( \frac{d\sigma^2}{1-k\sigma^2} + 
	\sigma^2\left[ d\theta^2 + 
	\sin^2\theta d\phi^2 \right] \right)
\end{equation}
\end{verbatim}

\include{Chapters/Chapter3}
\include{Chapters/Chapter4}
\include{Chapters/Chapter5}

%----------------------------------------------------------------------------------------
% Apéndices
%----------------------------------------------------------------------------------------

\appendix

% Incluir apéndices desde archivos separados si es necesario
%\include{Appendices/AppendixA}

%----------------------------------------------------------------------------------------
% Bibliografía
%----------------------------------------------------------------------------------------

\renewcommand{\bibname}{Bibliografía} % Para asegurarte de que el título sea correcto
\phantomsection % Necesario para que el enlace del marcador sea correcto

\printbibliography[heading=bibintoc]

\end{document}






